% Generated draft: requires mathematical and source review.
\documentclass[11pt]{article}
\usepackage[margin=25mm]{geometry}
\usepackage{fontspec}
\setmainfont{Latin Modern Roman}
\usepackage{amsmath,amssymb,mathtools,unicode-math}
\setmathfont{Latin Modern Math}
\usepackage{enumitem,needspace,xurl,hyperref}
\setlength{\parindent}{0pt}
\setlength{\parskip}{5pt}
\setlength{\emergencystretch}{3em}
\newcommand{\N}{\mathbb{N}}
\newcommand{\Z}{\mathbb{Z}}
\newcommand{\Q}{\mathbb{Q}}
\newcommand{\R}{\mathbb{R}}
\newcommand{\C}{\mathbb{C}}
\newcommand{\F}{\mathbb{F}}
\newcommand{\A}{\mathbb{A}}
\newcommand{\PP}{\mathbb{P}}
\newcommand{\E}{\mathbb{E}}
\newcommand{\eps}{\varepsilon}
\newcommand{\dd}{\,\mathrm{d}}
\begin{document}
\section*{UnsolvedMath ID: 20002063}
\textbf{Source problem code:} catalogue entry
\subsection{Definitions and mathematical statement}
Ekholm's Question 1.4 has two tentative parts: (a) compare linearized contact homology with generating-function homology for a possibly compact Legendrian \(L\subset J^1(\R^n)\), with the generating-function definition left in brackets; (b) ask whether a generating function exists if and only if linearized contact homology is ``(uniquely) defined.'' Part (b) literally prints \(L\subset\R^n\), which does not specify a contact manifold or Legendrian dimension. The supplied source report says this is genuine source wording, not an OCR error; the likely \(J^1(\R^n)\) interpretation remains explicit rather than silent.

The aid's precise partial answer concerns closed chord-generic Legendrian knots
\[
 L\subset J^1(\R)=T^*\R\times\R_z,
 \qquad\alpha=\dd z-p\,\dd x,
\]
with an integer Maslov grading, in particular the rotation-zero grading regime, and coefficients \(\F_2\). Let \((\mathcal A_L,\partial)\) be the Chekanov--Eliashberg DGA. A graded augmentation is a unital graded DGA map
\(\epsilon:(\mathcal A_L,\partial)\to(\F_2,0)\).
Linearizing after this augmentation gives a chain complex generated by Reeb chords and homology \(LCH_*^\epsilon(L)\). The augmentation is part of the datum, not automatically unique.

An admissible generating family \(F:\R_x\times\R_\eta^N\to\R\) has regular fibre-critical locus \(\{\partial_\eta F=0\}\) parametrizing \(L\) by
\[
 (x,\eta)\longmapsto(x,\partial_xF,F).
\]
Use the linear-at-infinity class stipulated by the aid, with its required tameness/stabilization conventions. For the difference function
\(\delta_F(x,\eta,\eta')=F(x,\eta')-F(x,\eta)\), a standard grading convention is
\[
 GH_k(F)=H_{k+N+1}(\delta_F^{\leq\Omega},
                         \delta_F^{\leq\varepsilon};\F_2),
\]
where \(\varepsilon>0\) is below all positive critical values and \(\Omega\) is above them. The exact class at infinity and grading conventions must be matched to the intended comparison; the original brackets do not define them.

In this knot setting the aid reports the equivalence
\[
 L\text{ admits an admissible generating family}
 \quad\Longleftrightarrow\quad
 L\text{ has at least one graded augmentation}.
\]
It attributes the route to normal-ruling equivalences. Furthermore, for each such family \(F\) it reports a compatible augmentation \(\epsilon_F\) with
\[
 GH_*(F)\simeq LCH_*^{\epsilon_F}(L)
\]
under matched grading conventions. This has the quantifier \(\forall F\ \exists\epsilon_F\); it is not an assertion that every augmentation and every family give the same invariant, nor a bijection or uniqueness theorem.

The aid reports the Melvin--Shrestha Legendrian representative of the mirror knot \(m(8_{21})\) as a counterexample to the uniqueness reading. It admits a generating family but has graded augmentations \(\epsilon_0,\epsilon_1\) with distinct polynomials
\[
 P_{\epsilon_i}(t)=\sum_k
 \dim_{\F_2}LCH_k^{\epsilon_i}(L)\,t^k,
 \qquad P_{\epsilon_0}\ne P_{\epsilon_1}.
\]
The front diagram, DGA and numerical polynomials are not supplied, so this example is preserved as a named literature report rather than reconstructed data. It refutes uniqueness of augmentation-dependent graded homology, not the existence equivalence. The unrestricted higher-dimensional equality in (a) is not resolved here. The source's reported evidence for a Legendrian formed as the double of a compact \(n\)-manifold with boundary \(M\subset\R^n\) lacks its construction and comparison hypotheses; the trailing 4 is a footnote marker, not an exponent.

Theme 1.5 is adjoining research context, explicitly not a question: it mentions Fukaya--Oh/Witten degeneration of holomorphic discs to gradient trees; work of Lenny Ng and Zhu Ke on relative contact homology; Mikhalkin's amoeba/fibre-shrinking perspective in a Lagrangian fibration; and comparisons involving the Chas--Sullivan loop-space product and Floer cup product. None of these remarks supplies a further theorem or resolves the missing hypotheses of Question 1.4.
\subsection{Short English statement}
In the stated knot, grading and coefficient setting, generating families correspond to existence of an augmentation and select a matching linearized homology. They do not force uniqueness: the aid reports a generating-family knot with two different linearized homologies. Higher-dimensional and function-space qualifications remain open in this record.
\subsection{Sources}
\begin{itemize}
\item Permanent UnsolvedMath ID \(20002063\), exact public code catalogue entry; local registry \texttt{aim-geometry-notes.json}, \texttt{notes[400]}, canonical tag question.
\item AIM workshop ``Holomorphic curves in contact geometry,'' Question 1.4 attributed to Ekholm, followed by Theme 1.5: \url{https://aimath.org/WWN/contactgeom2/contactgeom2.pdf}. No PDF was accessed. The supplied verification report fixes the \(J^1\) coordinate convention, confirms the literal \(\R^n\) wording in (b), and identifies the footnote marker.
\item The aid names Melvin--Shrestha and their \(m(8_{21})\) computation, and invokes normal-ruling equivalences without full references. Theme attributions are Fukaya--Oh, Witten, Lenny Ng, Zhu Ke, Mikhalkin and Chas--Sullivan. Original literature and year metadata are absent.
\item Research attempt 1, ``Generating families, augmentations, and the missing quantifiers,'' labelled partial\_result, partially\_solved, proved and low novelty confidence; statement corrected\_verified. The clean field has additional theme-number corruption, so the original Question 1.4 identity is retained.
\end{itemize}
\end{document}
